\chapter{Literature review}\label{chap:lit}

This chapter places the thesis in four lines of work:
\begin{itemize}
    \item hallucination in vision-language models and its benchmarks;
    \item supervised span-level detectors and the synthetic data used to train them;
    \item uncertainty signals, both black-box and white-box;
    \item the evaluation methodology for comparing many configurations on one validation set.
\end{itemize}
Each section ends with the gap that this thesis addresses.

\section{Hallucination in vision-language models}

In language generation, a hallucination is content that is fluent but not supported by the source or by world knowledge~\cite{ji2023survey}. In vision-language models the source is the image, and early work measured object hallucination in captions: CHAIR counts mentioned objects that are absent from the image annotation~\cite{rohrbach2018object}, and POPE asks the model yes/no questions about the presence of objects~\cite{li2023pope}. Both work at the level of whole objects or whole answers. Later analyses linked hallucination to how the model attends to the image. OPERA, for example, shows that VLMs over-trust a few ``summary'' tokens and stop attending to visual evidence as the answer grows~\cite{huang2024opera}. This observation motivates the attention look-back ratio used in our probe (Section~\ref{sec:probe}).

Span-level benchmarks ask a different question: which characters of an answer are unsupported. The SemEval-2025 Mu-SHROOM task defined it for text-only LLM answers in fourteen languages, with soft gold labels obtained by aggregating several annotators~\cite{vazquez-etal-2025-semeval}. SHROOM-Visions 2026 brought the same design to images~\cite{vazquez2026shroomvisionsoverview}. Its data was written and checked by humans to obtain realistic, fine-grained hallucinations~\cite{mickus2026humansdreamelectricsheep}. It adds five categories (invention, mischaracterization, OCR, miscounting and other) and a label-aware metric, Cor+Lbl. Chapter~\ref{chap:problem} defines the task formally.

\textbf{Gap.} Object- and answer-level benchmarks cannot evaluate span-level, category-aware, calibrated output. SHROOM-Visions is the first benchmark that does, so the methods that work for it are not yet established.

\section{Supervised detectors and synthetic training data}

The strongest Mu-SHROOM systems fine-tuned LLMs to tag hallucinated spans, often on synthetic or weakly supervised data~\cite{rykov-etal-2025-smurfcat,mo-etal-2025-team,kobus-etal-2025-atlantis}, or prompted large models with translation strategies for low-resource languages~\cite{karkani-etal-2025-ails}. Parameter-efficient fine-tuning makes this practical on a single GPU. LoRA learns low-rank updates of frozen weights~\cite{hu2022lora}, and QLoRA combines it with 4-bit quantisation of the frozen model~\cite{dettmers2023qlora}. In SHROOM-Visions, the best system fine-tuned five VLM judges from four model families and combined them by majority vote, which ranked first in three of the four languages~\cite{vroomvroom2026}. Our own shared-task entry fine-tuned Qwen3-VL-8B~\cite{bai2025qwen3} with upsampled and synthetically injected training data~\cite{rattanajaratkul2026smurfcat}.

Synthetic data is attractive because annotated hallucinations are expensive and rare categories are very rare: in SHROOM-Visions, only 6\% of training answers contain an OCR error. Synthetic hallucinations are typically made by perturbing real answers~\cite{zhang2024perturbation} or by prompting an LLM under controlled instructions~\cite{xie2024controlled}. Both papers note a risk. If synthetic spans differ from real ones in surface form, a detector can learn the construction artifact instead of the grounding error. Our earlier analysis found exactly this for template-based injections. A text-only classifier separates them from real spans with an AUC of 0.92--1.0 (Section~\ref{sec:exp-synth}).

\textbf{Gap.} Synthetic data is usually added together with other changes, such as upsampling, more training steps or a different label distribution, so its own effect is rarely known. No study we are aware of tests synthetic construction methods one at a time against the same baseline for span-level VLM hallucination detection.

\section{Uncertainty signals}

\paragraph{Black-box consistency.}
If a model knows a fact, repeated samples tend to agree. If it hallucinates, they tend to diverge. SelfCheckGPT turns this into a detector by checking each sentence of an answer against sampled answers~\cite{manakul2023selfcheckgpt}. Semantic entropy clusters samples by meaning before computing entropy~\cite{kuhn2023semantic,farquhar2024semantic}. Self-consistency decoding takes a majority vote over sampled reasoning paths to improve accuracy~\cite{wang2023selfconsistency}. For Mu-SHROOM, a training-free system localised spans by matching sliding windows of the answer against sampled generations~\cite{vemula2025keepitsimple}. These methods sample the \emph{generator}, the model whose answer is being checked. In SHROOM-Visions the generator is not available, so we apply the idea to the \emph{detector} instead. We sample the fine-tuned judge several times and vote on characters.

\paragraph{White-box probes.}
Internal states of a model also carry information about hallucination. INSIDE uses the eigenvalues of the covariance of hidden states over sampled answers~\cite{chen2024inside}. Lookback Lens classifies spans from the ratio of attention to the context versus to newly generated tokens, and matches classifiers that use full hidden states~\cite{chuang2024lookback}. Layer-wise analysis shows that signal accumulated across layers beats the last layer alone~\cite{kim2025layerwise}. For span output, recent work treats hallucination detection as sequence labelling and decodes it with a conditional random field~\cite{lafferty2001crf,itkin2026temporal}. This follows the BiLSTM--CRF line of named-entity recognisers~\cite{huang2015bilstmcrf,lample2016neural}. A CRF learns span contiguity directly, so no per-token threshold has to be searched on validation.

\paragraph{Combining signals.}
Ensembles reduce error when their members err differently~\cite{dietterich2000ensemble}. The top SHROOM-Visions system reports that adding internal-activation probes to its judge committee did not help, and attributes this to errors that were highly correlated with the judges~\cite{vroomvroom2026}.

\textbf{Gap.} Consistency methods have been used for answer-level verdicts and for text-only span detection. They have not been used for multi-label character spans with soft probabilities. Probes have been evaluated alone or as committee members without much success. Whether a weak probe can still help a self-consistency vote, and how strongly its errors correlate with the vote, is an open question.

\section{Reliable evaluation under validation reuse}

When many configurations are compared on the same holdout set, the best one's score overstates its true performance~\cite{dwork2015reusable}. The bias grows with the number of comparisons and with their adaptivity, for example per-language or per-category choices made one cell at a time. Shared-task systems are especially exposed: validation is small, many variants are tried, and test is used only once at the end. Our own earlier per-cell selection among uncertainty signals is an example. It looked strong on validation and dropped sharply on test~\cite{rattanajaratkul2026smurfcat}. Standard remedies are to limit the number of adaptive comparisons, to confirm a choice on data that did not take part in it, and to report variance over training seeds and resampling~\cite{efron1993bootstrap}.

\textbf{Gap.} Span-level hallucination detection results are rarely reported with seed variance, confirmation data or a fixed search budget, so it is often unclear which component of a system helps.

\section{Position of this thesis}

The thesis connects these lines in one architecture and tests each part in isolation:
\begin{itemize}
    \item \textbf{Judge.} A supervised generative judge is chosen among four VLM families under a fixed recipe.
    \item \textbf{Training data.} Synthetic data constructions, including a new retrieval-grounded one, are tested one at a time against the same baseline.
    \item \textbf{Uncertainty.} A self-consistency vote over the judge's own samples, adapted to multi-label character spans, and a multi-label CRF probe on the frozen judge are combined and analysed for error correlation.
\end{itemize}
Every comparison follows a pre-registered protocol designed against validation reuse.
